% =====================================================================
%  Appendix B: Algebra, Trigonometry, and Complex Numbers Quick Reference
% =====================================================================

\chapter{Algebra, Trigonometry, and Complex Numbers Quick Reference}
\label{app:algebra-trig-complex}

This appendix collects the algebra, trigonometry, and complex-number facts used in the
book.  It is meant to be scanned.

\[
\begin{array}{c|c}
\text{Tool} & \text{Where it appears} \\ \hline
\text{factoring} & \text{limits, domains, singularities}\\[1mm]
\text{completing the square} & \text{circles, ellipses, quadratic forms}\\[1mm]
\text{inequalities} & \text{estimates, domains, convergence}\\[1mm]
\text{trigonometry} & \text{polar, cylindrical, spherical coordinates}\\[1mm]
\text{complex numbers} & \text{plane rotations and roots}\\[1mm]
\text{Euler's formula} & \text{oscillation, rotation, roots of unity}
\end{array}
\]

% ---------------------------------------------------------------------
\section{Algebraic manipulation and inequalities}
\label{app:B.1}

\subsection*{Factoring patterns}

\[
\boxed{
    a^2-b^2=(a-b)(a+b)
}
\]

\[
\boxed{
    a^2+2ab+b^2=(a+b)^2
}
\]

\[
\boxed{
    a^2-2ab+b^2=(a-b)^2
}
\]

\[
\boxed{
    a^3-b^3=(a-b)(a^2+ab+b^2)
}
\]

\[
\boxed{
    a^3+b^3=(a+b)(a^2-ab+b^2)
}
\]

\[
\boxed{
    x^2+(p+q)x+pq=(x+p)(x+q)
}
\]

\subsection*{Quadratic formula}

If
\[
    ax^2+bx+c=0,
    \qquad
    a\ne0,
\]
then
\[
\boxed{
    x=\frac{-b\pm\sqrt{b^2-4ac}}{2a}.
}
\]

The discriminant is
\[
\boxed{
    \Delta=b^2-4ac.
}
\]

\[
\begin{array}{c|c}
\Delta>0 & \text{two distinct real roots}\\
\Delta=0 & \text{one repeated real root}\\
\Delta<0 & \text{two complex conjugate roots}
\end{array}
\]

\subsection*{Completing the square}

\[
\boxed{
    x^2+bx
    =
    \left(x+\frac b2\right)^2-\left(\frac b2\right)^2.
}
\]

More generally,
\[
\boxed{
    ax^2+bx+c
    =
    a\left(x+\frac{b}{2a}\right)^2
    +
    c-\frac{b^2}{4a}
}
\]
for \(a\ne0\).

\subsection*{Circle form}

\[
\boxed{
    (x-h)^2+(y-k)^2=R^2
}
\]

Center:
\[
    (h,k).
\]

Radius:
\[
    R.
\]

\subsection*{Distance on the line}

\[
\boxed{
    |x-a|
}
\]
is the distance from \(x\) to \(a\).

\[
\boxed{
    |x-a|<r
    \quad\Longleftrightarrow\quad
    a-r<x<a+r.
}
\]

\[
\boxed{
    |x-a|\le r
    \quad\Longleftrightarrow\quad
    a-r\le x\le a+r.
}
\]

\[
\boxed{
    |x-a|>r
    \quad\Longleftrightarrow\quad
    x<a-r \text{ or } x>a+r.
}
\]

\subsection*{Absolute value rules}

\[
\boxed{
    |ab|=|a|\,|b|
}
\]

\[
\boxed{
    \left|\frac ab\right|=\frac{|a|}{|b|}
    \qquad (b\ne0)
}
\]

\[
\boxed{
    |a+b|\le |a|+|b|
}
\]
This is the triangle inequality.

\[
\boxed{
    \bigl||a|-|b|\bigr|\le |a-b|.
}
\]

\subsection*{Inequality rules}

Adding or subtracting the same number preserves the inequality:
\[
    a<b
    \quad\Longrightarrow\quad
    a+c<b+c.
\]

Multiplying by a positive number preserves the inequality:
\[
    a<b,\ c>0
    \quad\Longrightarrow\quad
    ac<bc.
\]

Multiplying by a negative number reverses the inequality:
\[
    a<b,\ c<0
    \quad\Longrightarrow\quad
    ac>bc.
\]

\subsection*{Quadratic inequalities}

To solve
\[
    ax^2+bx+c>0
\]
or
\[
    ax^2+bx+c<0,
\]
use this process:

\[
\begin{array}{ll}
1. & \text{Find the zeros.}\\
2. & \text{Put the zeros on the number line.}\\
3. & \text{Test one point in each interval.}\\
4. & \text{Include endpoints only for }\le\text{ or }\ge.
\end{array}
\]

\subsection*{Exponent rules}

\[
\boxed{
    a^m a^n=a^{m+n}
}
\]

\[
\boxed{
    \frac{a^m}{a^n}=a^{m-n}
    \qquad (a\ne0)
}
\]

\[
\boxed{
    (a^m)^n=a^{mn}
}
\]

\[
\boxed{
    (ab)^n=a^n b^n
}
\]

\[
\boxed{
    a^{-n}=\frac1{a^n}
    \qquad (a\ne0)
}
\]

\[
\boxed{
    a^{1/n}=\sqrt[n]{a}
}
\]

\subsection*{Logarithm rules}

For \(a>0\), \(b>0\), and \(r\in\mathbb R\),

\[
\boxed{
    \ln(ab)=\ln a+\ln b
}
\]

\[
\boxed{
    \ln\left(\frac ab\right)=\ln a-\ln b
}
\]

\[
\boxed{
    \ln(a^r)=r\ln a
}
\]

\[
\boxed{
    e^{\ln a}=a
}
\]

\[
\boxed{
    \ln(e^x)=x
}
\]

\subsection*{Common warnings}

Do not divide by a quantity unless it is known to be nonzero.

Do not take square roots of both sides of an inequality without tracking signs.

\[
    x^2>9
\]
means
\[
    x<-3
    \quad\text{or}\quad
    x>3,
\]
not just \(x>3\).

The identity
\[
    \sqrt{x^2}=|x|
\]
is usually safer than writing \(\sqrt{x^2}=x\).

% ---------------------------------------------------------------------
\section{Trigonometric identities}
\label{app:B.2}

\subsection*{Unit circle}

On the unit circle,
\[
\boxed{
    (\cos\theta,\sin\theta)
}
\]
is the point at angle \(\theta\).

\[
\boxed{
    \cos^2\theta+\sin^2\theta=1.
}
\]

\subsection*{Radians}

\[
\boxed{
    2\pi\text{ radians}=360^\circ.
}
\]

\[
\boxed{
    \pi\text{ radians}=180^\circ.
}
\]

\[
\boxed{
    s=r\theta
}
\]
where \(s\) is arc length, \(r\) is radius, and \(\theta\) is measured in radians.

\subsection*{Common angles}

\[
\begin{array}{c|cc}
\theta & \cos\theta & \sin\theta \\ \hline
0 & 1 & 0\\[1mm]
\dfrac{\pi}{6} & \dfrac{\sqrt3}{2} & \dfrac12\\[3mm]
\dfrac{\pi}{4} & \dfrac{\sqrt2}{2} & \dfrac{\sqrt2}{2}\\[3mm]
\dfrac{\pi}{3} & \dfrac12 & \dfrac{\sqrt3}{2}\\[3mm]
\dfrac{\pi}{2} & 0 & 1\\[2mm]
\pi & -1 & 0\\[1mm]
\dfrac{3\pi}{2} & 0 & -1\\[3mm]
2\pi & 1 & 0
\end{array}
\]

\subsection*{Signs by quadrant}

\[
\begin{array}{c|cc}
\text{Quadrant} & \cos\theta & \sin\theta \\ \hline
\text{I} & + & +\\
\text{II} & - & +\\
\text{III} & - & -\\
\text{IV} & + & -
\end{array}
\]

\subsection*{Other trigonometric functions}

\[
\boxed{
    \tan\theta=\frac{\sin\theta}{\cos\theta}
}
\]

\[
\boxed{
    \sec\theta=\frac1{\cos\theta}
}
\]

\[
\boxed{
    \csc\theta=\frac1{\sin\theta}
}
\]

\[
\boxed{
    \cot\theta=\frac{\cos\theta}{\sin\theta}
}
\]

\subsection*{Pythagorean identities}

\[
\boxed{
    \sin^2\theta+\cos^2\theta=1
}
\]

\[
\boxed{
    1+\tan^2\theta=\sec^2\theta
}
\]

\[
\boxed{
    1+\cot^2\theta=\csc^2\theta
}
\]

\subsection*{Even and odd identities}

\[
\boxed{
    \cos(-\theta)=\cos\theta
}
\]

\[
\boxed{
    \sin(-\theta)=-\sin\theta
}
\]

\[
\boxed{
    \tan(-\theta)=-\tan\theta
}
\]

\subsection*{Periodicity}

\[
\boxed{
    \sin(\theta+2\pi)=\sin\theta
}
\]

\[
\boxed{
    \cos(\theta+2\pi)=\cos\theta
}
\]

\[
\boxed{
    \tan(\theta+\pi)=\tan\theta
}
\]

\subsection*{Angle addition formulas}

\[
\boxed{
    \cos(\alpha+\beta)
    =
    \cos\alpha\cos\beta-\sin\alpha\sin\beta
}
\]

\[
\boxed{
    \sin(\alpha+\beta)
    =
    \sin\alpha\cos\beta+\cos\alpha\sin\beta
}
\]

\[
\boxed{
    \cos(\alpha-\beta)
    =
    \cos\alpha\cos\beta+\sin\alpha\sin\beta
}
\]

\[
\boxed{
    \sin(\alpha-\beta)
    =
    \sin\alpha\cos\beta-\cos\alpha\sin\beta
}
\]

\subsection*{Double-angle formulas}

\[
\boxed{
    \sin(2\theta)=2\sin\theta\cos\theta
}
\]

\[
\boxed{
    \cos(2\theta)=\cos^2\theta-\sin^2\theta
}
\]

\[
\boxed{
    \cos(2\theta)=2\cos^2\theta-1
}
\]

\[
\boxed{
    \cos(2\theta)=1-2\sin^2\theta
}
\]

\subsection*{Half-angle identities}

\[
\boxed{
    \cos^2\theta=\frac{1+\cos(2\theta)}{2}
}
\]

\[
\boxed{
    \sin^2\theta=\frac{1-\cos(2\theta)}{2}
}
\]

\subsection*{Useful shifts}

\[
\boxed{
    \cos(\theta+\pi)=-\cos\theta
}
\]

\[
\boxed{
    \sin(\theta+\pi)=-\sin\theta
}
\]

\[
\boxed{
    \cos\left(\theta+\frac{\pi}{2}\right)=-\sin\theta
}
\]

\[
\boxed{
    \sin\left(\theta+\frac{\pi}{2}\right)=\cos\theta
}
\]

\subsection*{Inverse trig derivatives}

\[
\boxed{
    \frac{d}{dx}\sin^{-1}x=\frac1{\sqrt{1-x^2}}
}
\]

\[
\boxed{
    \frac{d}{dx}\cos^{-1}x=-\frac1{\sqrt{1-x^2}}
}
\]

\[
\boxed{
    \frac{d}{dx}\tan^{-1}x=\frac1{1+x^2}
}
\]

\subsection*{Polar conversion}

\[
\boxed{
    x=r\cos\theta,
    \qquad
    y=r\sin\theta.
}
\]

\[
\boxed{
    r^2=x^2+y^2.
}
\]

\[
\boxed{
    \tan\theta=\frac yx
}
\]
with quadrant care.

\subsection*{Common warnings}

Radians are the default in calculus.

\[
    \sin^2\theta
\]
means
\[
    (\sin\theta)^2,
\]
not
\[
    \sin(\theta^2).
\]

The equation
\[
    \tan\theta=\frac yx
\]
does not determine the quadrant by itself.

% ---------------------------------------------------------------------
\section{Complex numbers and rotations}
\label{app:B.3}

\subsection*{Complex numbers}

\[
\boxed{
    z=a+bi
}
\]
where
\[
\boxed{
    i^2=-1.
}
\]

Real part:
\[
\boxed{
    \operatorname{Re}(a+bi)=a.
}
\]

Imaginary part:
\[
\boxed{
    \operatorname{Im}(a+bi)=b.
}
\]

Plane point:
\[
\boxed{
    a+bi
    \longleftrightarrow
    (a,b).
}
\]

\subsection*{Addition and subtraction}

\[
\boxed{
    (a+bi)+(c+di)=(a+c)+(b+d)i.
}
\]

\[
\boxed{
    (a+bi)-(c+di)=(a-c)+(b-d)i.
}
\]

\subsection*{Multiplication}

\[
\boxed{
    (a+bi)(c+di)=(ac-bd)+(ad+bc)i.
}
\]

Multiplication by \(i\):
\[
\boxed{
    i(a+bi)=-b+ai.
}
\]

As a plane rotation:
\[
\boxed{
    (a,b)\mapsto(-b,a).
}
\]
This is a \(90^\circ\) counterclockwise rotation.

\subsection*{Conjugate}

\[
\boxed{
    \overline{a+bi}=a-bi.
}
\]

Geometrically, conjugation reflects across the real axis:
\[
    (a,b)\mapsto(a,-b).
\]

\subsection*{Modulus}

\[
\boxed{
    |a+bi|=\sqrt{a^2+b^2}.
}
\]

\[
\boxed{
    z\overline z=|z|^2.
}
\]

\subsection*{Division}

To divide by
\[
    c+di,
\]
multiply top and bottom by
\[
    c-di.
\]

\[
\boxed{
    \frac{a+bi}{c+di}
    =
    \frac{(a+bi)(c-di)}{c^2+d^2}
}
\]
for
\[
    c+di\ne0.
\]

\subsection*{Polar form}

If
\[
    z=a+bi\ne0,
\]
then
\[
\boxed{
    z=r(\cos\theta+i\sin\theta)
}
\]
where
\[
\boxed{
    r=|z|=\sqrt{a^2+b^2}.
}
\]

The angle \(\theta\) satisfies
\[
    \cos\theta=\frac ar,
    \qquad
    \sin\theta=\frac br.
\]

\subsection*{Multiplication in polar form}

If
\[
    z=r(\cos\alpha+i\sin\alpha),
\]
and
\[
    w=s(\cos\beta+i\sin\beta),
\]
then
\[
\boxed{
    zw
    =
    rs\bigl(\cos(\alpha+\beta)+i\sin(\alpha+\beta)\bigr).
}
\]

Complex multiplication:

\[
\boxed{
    \text{multiplies lengths and adds angles.}
}
\]

\subsection*{De Moivre's formula}

\[
\boxed{
    \bigl(r(\cos\theta+i\sin\theta)\bigr)^n
    =
    r^n(\cos(n\theta)+i\sin(n\theta)).
}
\]

For unit complex numbers:
\[
\boxed{
    (\cos\theta+i\sin\theta)^n
    =
    \cos(n\theta)+i\sin(n\theta).
}
\]

\subsection*{Common warnings}

The number \(0\) has no polar angle.

Angles in polar form are not unique:
\[
    \theta
    \quad\text{and}\quad
    \theta+2\pi k
\]
describe the same direction.

Complex multiplication is not componentwise multiplication:
\[
    (a,b)(c,d)\ne(ac,bd).
\]

% ---------------------------------------------------------------------
\section{Euler's formula}
\label{app:B.4}

\subsection*{Euler form}

\[
\boxed{
    e^{i\theta}=\cos\theta+i\sin\theta.
}
\]

Therefore
\[
\boxed{
    |e^{i\theta}|=1.
}
\]

\subsection*{Polar form with Euler notation}

\[
\boxed{
    r(\cos\theta+i\sin\theta)=re^{i\theta}.
}
\]

So every nonzero complex number can be written as
\[
\boxed{
    z=re^{i\theta}.
}
\]

\subsection*{Multiplication}

\[
\boxed{
    (re^{i\alpha})(se^{i\beta})=rs e^{i(\alpha+\beta)}.
}
\]

\subsection*{Powers}

\[
\boxed{
    (re^{i\theta})^n=r^n e^{in\theta}.
}
\]

\subsection*{Rotation}

Multiplying by
\[
    e^{i\theta}
\]
rotates the plane by angle \(\theta\) counterclockwise and keeps lengths unchanged.

\[
\boxed{
    z\mapsto e^{i\theta}z
}
\]

\subsection*{Special values}

\[
\boxed{
    e^{i0}=1
}
\]

\[
\boxed{
    e^{i\pi/2}=i
}
\]

\[
\boxed{
    e^{i\pi}=-1
}
\]

\[
\boxed{
    e^{2\pi i}=1
}
\]

\[
\boxed{
    e^{i\pi}+1=0.
}
\]

\subsection*{Sine and cosine from exponentials}

\[
\boxed{
    \cos\theta=\frac{e^{i\theta}+e^{-i\theta}}{2}
}
\]

\[
\boxed{
    \sin\theta=\frac{e^{i\theta}-e^{-i\theta}}{2i}
}
\]

\subsection*{Euler formula and trig identities}

The identity
\[
    e^{i(\alpha+\beta)}=e^{i\alpha}e^{i\beta}
\]
contains the angle addition formulas:
\[
    \cos(\alpha+\beta)
    =
    \cos\alpha\cos\beta-\sin\alpha\sin\beta,
\]
\[
    \sin(\alpha+\beta)
    =
    \sin\alpha\cos\beta+\cos\alpha\sin\beta.
\]

\subsection*{Common warning}

The symbol
\[
    e^{i\theta}
\]
does not mean growth away from the origin.  It lies on the unit circle.  It represents
rotation.

% ---------------------------------------------------------------------
\section{Complex roots and geometry}
\label{app:B.5}

\subsection*{\(n\)th roots of a complex number}

Let
\[
    w=Re^{i\Theta},
    \qquad
    R>0.
\]
The solutions of
\[
    z^n=w
\]
are
\[
\boxed{
    z_k
    =
    R^{1/n}e^{i(\Theta+2\pi k)/n},
    \qquad
    k=0,1,\ldots,n-1.
}
\]

There are \(n\) distinct roots.

They lie on the circle of radius
\[
    R^{1/n}.
\]

Their angles are spaced by
\[
\boxed{
    \frac{2\pi}{n}.
}
\]

\subsection*{Roots of unity}

The solutions of
\[
    z^n=1
\]
are
\[
\boxed{
    z_k=e^{2\pi i k/n},
    \qquad
    k=0,1,\ldots,n-1.
}
\]

These are the \(n\)th roots of unity.

They form a regular \(n\)-gon on the unit circle.

\subsection*{Primitive root}

\[
\boxed{
    \omega=e^{2\pi i/n}
}
\]

Then the \(n\)th roots of unity are
\[
\boxed{
    1,\omega,\omega^2,\ldots,\omega^{n-1}.
}
\]

Also,
\[
\boxed{
    \omega^n=1.
}
\]

\subsection*{Sum of roots of unity}

For \(n>1\),
\[
\boxed{
    1+\omega+\omega^2+\cdots+\omega^{n-1}=0.
}
\]

Reason:
\[
    1+\omega+\cdots+\omega^{n-1}
    =
    \frac{1-\omega^n}{1-\omega}
    =
    0,
\]
because
\[
    \omega^n=1
    \qquad
    \text{and}
    \qquad
    \omega\ne1.
\]

\subsection*{Square roots}

To solve
\[
    z^2=Re^{i\Theta},
\]
use
\[
\boxed{
    z_0=\sqrt R e^{i\Theta/2},
    \qquad
    z_1=\sqrt R e^{i(\Theta+2\pi)/2}.
}
\]

The two roots are opposites.

\subsection*{Examples}

Fourth roots of unity:
\[
\boxed{
    1,\quad i,\quad -1,\quad -i.
}
\]

Cube roots of unity:
\[
\boxed{
    1,\quad
    -\frac12+\frac{\sqrt3}{2}i,\quad
    -\frac12-\frac{\sqrt3}{2}i.
}
\]

Solutions of
\[
    z^3=8:
\]
write
\[
    8=8e^{i0}.
\]
Then
\[
    z_k=2e^{2\pi i k/3},
    \qquad
    k=0,1,2.
\]
So
\[
\boxed{
    z=2,\quad
    -1+\sqrt3\,i,\quad
    -1-\sqrt3\,i.
}
\]

Solutions of
\[
    z^3=-8:
\]
write
\[
    -8=8e^{i\pi}.
\]
Then
\[
    z_k=2e^{i(\pi+2\pi k)/3},
    \qquad
    k=0,1,2.
\]
So
\[
\boxed{
    z=1+\sqrt3\,i,\quad
    -2,\quad
    1-\sqrt3\,i.
}
\]

\subsection*{Geometric meaning}

The map
\[
    z\mapsto z^n
\]
does two things:

\[
\boxed{
    \text{length } r \mapsto r^n
}
\]

\[
\boxed{
    \text{angle } \theta \mapsto n\theta.
}
\]

Taking \(n\)th roots reverses those operations:

\[
\boxed{
    \text{length } R \mapsto R^{1/n}
}
\]

\[
\boxed{
    \text{angle } \Theta \mapsto \frac{\Theta+2\pi k}{n}.
}
\]

The extra \(2\pi k\) is why there are \(n\) roots.

% ---------------------------------------------------------------------
\section*{One-page formula summary}

\[
\boxed{
    a^2-b^2=(a-b)(a+b)
}
\]

\[
\boxed{
    x=\frac{-b\pm\sqrt{b^2-4ac}}{2a}
}
\]

\[
\boxed{
    |a+b|\le |a|+|b|
}
\]

\[
\boxed{
    \sin^2\theta+\cos^2\theta=1
}
\]

\[
\boxed{
    \cos(\alpha+\beta)
    =
    \cos\alpha\cos\beta-\sin\alpha\sin\beta
}
\]

\[
\boxed{
    \sin(\alpha+\beta)
    =
    \sin\alpha\cos\beta+\cos\alpha\sin\beta
}
\]

\[
\boxed{
    x=r\cos\theta,\qquad y=r\sin\theta
}
\]

\[
\boxed{
    z=a+bi
}
\]

\[
\boxed{
    i^2=-1
}
\]

\[
\boxed{
    |a+bi|=\sqrt{a^2+b^2}
}
\]

\[
\boxed{
    \overline{a+bi}=a-bi
}
\]

\[
\boxed{
    z\overline z=|z|^2
}
\]

\[
\boxed{
    z=r(\cos\theta+i\sin\theta)=re^{i\theta}
}
\]

\[
\boxed{
    e^{i\theta}=\cos\theta+i\sin\theta
}
\]

\[
\boxed{
    (re^{i\theta})^n=r^n e^{in\theta}
}
\]

\[
\boxed{
    z^n=Re^{i\Theta}
    \quad\Longrightarrow\quad
    z_k=R^{1/n}e^{i(\Theta+2\pi k)/n}
}
\]

\[
\boxed{
    z^n=1
    \quad\Longrightarrow\quad
    z_k=e^{2\pi i k/n}
}
\]